% Pursuing Stacks (cotes 134-2 to 134-8): archive cote 134-4, site page 69 (PDF page 70), typescript page 157.
% typescript. Draft compared: ps3.tex:3557-3624.
% AI-transcribed by Claude Opus 5.5 (claude-opus-5-5, adaptive thinking, effort high), Message Batches API, batch ps2 (msgbatch_01NVwr1Yuo8qJf5rykeGBXPK), 2026-11-12; not yet checked by eye.
\scanpage{134-4/69}
\mspage{157}

Under these conditions,  the following conclusions hold:

a) $\underline{W} = f_{i}^{-1}(\underline{W})$ \quad (where $f_{i} = i_{\unsure{W|A}}i^{*}$ with yesterday's notations ) ,

b) The functor $\overline{f}_{i}$ from $\underline{W}^{-1}$(Cat) to itself induced by $f_{i}$ (which is
defined because of a)) is isomorphic (canonically) to the identi-%
ty functor, and hence is an equivalence.

The use we have for the \xout{diagram} \inserted[typed]{three} maps in the diagram (12) is complete-%
ly expressed in this lemma. The pretty obvious proof below would not
work at all if in a) above, we replace $\underline{W}$ in the first hand side by a $\underline{W}'$!
We have to prove that for a map $C \to C'$ in (Cat) , this is in $\underline{W}$ iff
$A_{/C} \to A_{/C'}$ is. Now this is seen from an obvious diagram chasing in the
diagram below, using saturation condition b) on $\underline{W}$ :
\diagram{\begin{tikzcd}
A_{/C} \ar[r]\ar[d] & A_{//C} \ar[d] & AxC \ar[l]\ar[r]\ar[d] & C \ar[d] \\
A_{/C'} \ar[r] & A_{//C'} & AxC' \ar[l]\ar[r] & C\hand{\unsure{'}}
\end{tikzcd}} ,

where all horizontal arrows are already known to be in $\underline{W}$ (the assumptions
in the lemma were designed for just that end). At the same time, we see
that the corresponding statements are equally true for the functors
$C \mapsto A_{//C}$ and $C \mapsto AxC$, and that \unsure{twe|two} consecutive among the four functors
we got from $H_{\underline{W}} = \underline{W}^{-1}$(Cat) to itself, deduced by localization by $\underline{W}$,
are canonically isomorphic, which proves b) by taking the composition
of the three isomorphims
\[ \hand{\gamma}(A_{/C}) \hand{\xrightarrow{\sim}} \hand{\gamma}(A_{//C}) \hand{\xrightarrow{\sim}} \hand{\gamma}(AxC) \hand{\xrightarrow{\sim}} \hand{\gamma}(C) \quad , \]
where $\hand{\gamma}$: (Cat)$\to H_{\underline{W}} = \underline{W}^{-1}$(Cat) is the canonical functor.

With this lemma, we have everything needed in order to write down
the full closed chain of implications, between various conditions on
($\underline{W}$,A,i) , from which to read off the ``key result'' we're after. Before
doing so, I would like still to make some preliminary comments on the
role of $\underline{W}$ , and \inserted{on} the nature of the conditions we have been led so far to
impose upon $\underline{W}$ .

A first feature that is striking, is that all conditions needed
are in the nature either of stability conditions (if such and such maps
are in $\underline{W}$ , so are others deduced from it), or conditions stating that
such and such unqualified maps (\xout{xxx}\inserted[typed]{the}projection $\hand{\Delta}_{1}$xC $\to$ C for any C,
say) are in $\underline{W}$ . We did not have any use of the only condition stated
so far, namely b) (if f $\hand{\in}$ $\underline{W}$, $\hand{\pi}_{o}$(f) is bijective) \deleted{which is} of a \uline{restric-}%
\uline{tive} type on the kind of arrows allowed in $\underline{W}$ - which is quite contrary
to my expectations. Thus, all conditions are trivially met if we take

\margin{69}
